\section{Resultados}

\subsection{Evolución de las disparidades departamentales}

La Tabla \ref{tab:estadisticos_descriptivos} presenta los principales
estadísticos descriptivos del Índice de Desarrollo Humano y del logaritmo del
PBI real per cápita relativo. El IDH promedio departamental aumentó de 0.3488
en 2007 a 0.5175 en 2019. La mediana también se desplazó de 0.3424 a 0.5114,
lo que confirma una mejora generalizada de los niveles de desarrollo humano
durante el periodo analizado.

El incremento del promedio estuvo acompañado por un aumento de la desviación
estándar, que pasó de 0.0693 a 0.0840. No obstante, el coeficiente de variación
disminuyó de 0.1986 a 0.1624. Esto indica que, en relación con el nivel medio
alcanzado, la dispersión del IDH se redujo. Por tanto, existe evidencia
descriptiva compatible con convergencia sigma relativa, aunque esta conclusión
debe contrastarse posteriormente con la regresión de convergencia beta.

En el caso del logaritmo del PBI real per cápita relativo, la media pasó de
$-0.2144$ en 2007 a $-0.1211$ en 2019, mientras que la mediana aumentó de
$-0.3496$ a $-0.2253$. La desviación estándar disminuyó de 0.6193 a 0.4750,
lo cual representa una reducción aproximada de 23.3\% en la dispersión
absoluta del indicador. El coeficiente de variación no se interpreta para esta
variable porque su media es negativa y cercana a cero.

Los estadísticos de Jarque--Bera presentan valores de probabilidad superiores
a 0.05 en todas las variables. En consecuencia, no se rechaza la hipótesis de
normalidad al nivel de significancia del 5\%. Esta evidencia respalda el empleo
de métodos paramétricos como primera aproximación, aunque posteriormente se
incorporan estimaciones no paramétricas para evaluar posibles relaciones no
lineales.

\begin{table}[H]
\centering
\caption{Estadísticos descriptivos de los indicadores departamentales}
\label{tab:estadisticos_descriptivos}
\small
\resizebox{\textwidth}{!}{
\begin{tabular}{lrrrrrrrrr}
\toprule
\textbf{Variable} &
\textbf{Media} &
\textbf{Mediana} &
\textbf{Máximo} &
\textbf{Mínimo} &
\textbf{D.E.} &
\textbf{C.V.} &
\textbf{Asimetría} &
\textbf{Curtosis} &
\textbf{$p$-valor JB} \\
\midrule
IDH 2007
& 0.3488 & 0.3424 & 0.4901 & 0.2147 & 0.0693 & 0.1986
& 0.2254 & 2.2817 & 0.6979 \\

IDH 2019
& 0.5175 & 0.5114 & 0.7073 & 0.3838 & 0.0840 & 0.1624
& 0.5220 & 2.5548 & 0.5251 \\

$\ln$ PBIpc relativo 2007
& $-0.2144$ & $-0.3496$ & 1.4650 & $-0.9515$ & 0.6193
& --- & 0.8269 & 3.2388 & 0.2476 \\

$\ln$ PBIpc relativo 2019
& $-0.1211$ & $-0.2253$ & 1.1337 & $-0.7302$ & 0.4750
& --- & 0.7989 & 3.0146 & 0.2790 \\
\bottomrule
\end{tabular}
}
\begin{flushleft}
\footnotesize
\textit{Nota.} D.E. = desviación estándar; C.V. = coeficiente de variación;
JB = prueba de Jarque--Bera. El coeficiente de variación no se presenta para
el logaritmo del PBI per cápita relativo porque su media es negativa.
\end{flushleft}
\end{table}

La evolución distributiva del IDH también puede observarse en la Figura
\ref{fig:densidad_idh}. La densidad estimada para 2019 se encuentra desplazada
hacia la derecha respecto de la distribución de 2007, reflejando una mejora
general de los niveles departamentales. Sin embargo, el desplazamiento no
implica que todos los departamentos hayan avanzado a la misma velocidad.

\begin{figure}[H]
\centering
\includegraphics[width=0.88\textwidth]{figura_1_densidad_idh.png}
\caption{Distribución kernel del Índice de Desarrollo Humano, 2007 y 2019}
\label{fig:densidad_idh}
\begin{flushleft}
\footnotesize
\textit{Nota.} Las curvas representan densidades kernel departamentales.
El desplazamiento hacia la derecha indica una mejora general de los niveles
de desarrollo humano.
\end{flushleft}
\end{figure}

La Figura \ref{fig:mapas_idh} muestra la distribución territorial del IDH
inicial y de su cambio acumulado. Los niveles iniciales evidencian diferencias
persistentes entre los departamentos, mientras que el mapa de variaciones
permite observar que los avances no se distribuyeron de manera uniforme. La
existencia de mejoras generalizadas no elimina, por tanto, las diferencias en
la intensidad con la que cada territorio amplió sus capacidades.

\begin{figure}[H]
\centering
\includegraphics[width=0.94\textwidth]{figura_2_mapas_idh.png}
\caption{Distribución departamental del IDH inicial y su cambio acumulado}
\label{fig:mapas_idh}
\begin{flushleft}
\footnotesize
\textit{Nota.} El panel inicial representa los niveles departamentales de
2007; el segundo muestra la variación acumulada entre 2007 y 2019.
\end{flushleft}
\end{figure}

\subsection{Autocorrelación espacial global}

Los resultados del índice global de Moran se presentan en la Tabla
\ref{tab:moran_global}. Ninguno de los estadísticos estimados resulta
significativo al 5\%. Para la matriz de banda de distancia inversa, el Moran
del IDH inicial es $-0.0273$ y el correspondiente a su cambio acumulado es
$-0.0662$. En ambos casos, los valores de probabilidad son ampliamente
superiores a 0.05.

En el PBI per cápita relativo inicial, el Moran alcanza 0.0692, pero tampoco es
estadísticamente significativo. El cambio económico presenta un valor de
$-0.0769$, igualmente no significativo. Por tanto, no se identifica un patrón
global estable mediante el cual departamentos con valores similares se
encuentren sistemáticamente próximos entre sí.

Los resultados del IDH se mantienen cuando se utilizan las matrices de
contigüidad queen y cuatro vecinos más próximos. Esta estabilidad sugiere que
la ausencia de autocorrelación global no depende exclusivamente de la matriz
principal. No obstante, la inexistencia de asociación global significativa no
descarta la presencia de relaciones locales o de heterogeneidad en los
coeficientes regionales.

\begin{table}[H]
\centering
\caption{Índice global de Moran}
\label{tab:moran_global}
\small
\begin{tabular}{llrr}
\toprule
\textbf{Variable} &
\textbf{Matriz} &
\textbf{Moran $I$} &
\textbf{$p$-valor} \\
\midrule
IDH 2007               & Banda inversa & $-0.0273$ & 0.4443 \\
Cambio del IDH          & Banda inversa & $-0.0662$ & 0.5796 \\
$\ln$ PBIpc 2007        & Banda inversa & 0.0692    & 0.1581 \\
Cambio de $\ln$ PBIpc   & Banda inversa & $-0.0769$ & 0.6265 \\
IDH 2007                & Queen         & $-0.0690$ & 0.5752 \\
Cambio del IDH          & Queen         & $-0.0277$ & 0.4525 \\
IDH 2007                & KNN, $k=4$    & $-0.0713$ & 0.5903 \\
Cambio del IDH          & KNN, $k=4$    & $-0.0799$ & 0.6194 \\
\bottomrule
\end{tabular}
\begin{flushleft}
\footnotesize
\textit{Nota.} La hipótesis nula corresponde a ausencia de autocorrelación
espacial global. Ningún resultado es significativo al 5\%.
\end{flushleft}
\end{table}

El diagrama de dispersión espacial de la Figura \ref{fig:moran_idh} confirma
que los departamentos se distribuyen entre los distintos cuadrantes sin
formar un agrupamiento global claramente dominante. La pendiente de la recta
corresponde al índice de Moran y se mantiene próxima a cero.

\begin{figure}[H]
\centering
\includegraphics[width=0.90\textwidth]{figura_3_moran_scatter.png}
\caption{Diagramas de dispersión de Moran para el IDH}
\label{fig:moran_idh}
\begin{flushleft}
\footnotesize
\textit{Nota.} Los ejes representan el valor estandarizado del indicador y
su rezago espacial. La pendiente corresponde al índice global de Moran.
\end{flushleft}
\end{figure}

\subsection{Convergencia beta global}

Los resultados de las regresiones globales evidencian comportamientos
diferenciados entre desarrollo humano e ingreso regional. En la estimación
OLS del IDH, el coeficiente asociado con el nivel inicial es positivo y
estadísticamente significativo:

\begin{equation}
\widehat{\Delta IDH_i}
=
0.1115
+
0.1637 IDH_{i,2007}.
\end{equation}

El coeficiente de 0.1637 presenta un valor de probabilidad de 0.0355. Esto
significa que los departamentos con mayor IDH inicial registraron, en promedio,
incrementos superiores durante 2007--2019. Por consiguiente, no se encuentra
evidencia de convergencia beta del desarrollo humano. El resultado es más
compatible con una dinámica de divergencia beta o de persistencia acumulativa,
a pesar de que el coeficiente de variación disminuyó.

El modelo OLS del IDH alcanza un $R^2$ de 0.1858 y un $R^2$ ajustado de
0.1488. El nivel explicativo es moderado, lo que indica que una proporción
importante de las variaciones departamentales responde a factores no incluidos
en una regresión basada únicamente en el nivel inicial.

En contraste, el coeficiente del PBI real per cápita relativo inicial es
negativo y altamente significativo:

\begin{equation}
\widehat{\Delta x_i}
=
0.0229
-
0.3282 x_{i,2007}.
\end{equation}

El coeficiente de $-0.3282$ presenta un valor de probabilidad inferior a
0.001. Esto implica que los departamentos con una posición económica inicial
más rezagada mejoraron relativamente más que los departamentos inicialmente
adelantados. Por tanto, se identifica convergencia beta económica durante
2007--2019.

El modelo económico alcanza un $R^2$ de 0.4404 y un $R^2$ ajustado de 0.4150.
La velocidad anual implícita de convergencia, calculada a partir del coeficiente
OLS, es aproximadamente:

\begin{equation}
v
=
-\frac{\ln(1-0.3282)}{12}
=
0.0331.
\end{equation}

Por tanto, la velocidad estimada de convergencia económica es cercana al
3.3\% anual. Este resultado debe interpretarse como una aproximación promedio
para el conjunto departamental y no como una tasa uniforme para todos los
territorios.

\subsection{Modelos espaciales y robustez}

La Tabla \ref{tab:modelos_regresion} compara las estimaciones OLS, SEM y
Durbin-SEM. En el caso del IDH, el coeficiente inicial permanece positivo y
significativo bajo todas las especificaciones. Con la matriz de banda inversa,
el SEM estima un coeficiente de 0.1712, mientras que con la matriz queen el
valor asciende a 0.1882. Los resultados confirman que la ausencia de
convergencia beta del IDH no depende del modelo elegido.

El parámetro asociado con el rezago espacial del IDH inicial no es
significativo en la matriz de distancia. En la especificación queen alcanza
0.2782 y presenta un valor de probabilidad de 0.0710, por lo que solo sería
significativo bajo un criterio del 10\%. El AIC más bajo continúa
correspondiendo al OLS, aunque la diferencia respecto del Durbin-SEM queen es
reducida.

Para el PBI per cápita relativo, todos los modelos mantienen un coeficiente
inicial negativo y altamente significativo. En el SEM con banda inversa, el
coeficiente alcanza $-0.3666$, mientras que en el SEM queen asciende en valor
absoluto a $-0.3792$. El modelo queen presenta el menor AIC del conjunto
económico, con un valor de $-1.3505$.

Los coeficientes de los rezagos espaciales del nivel inicial no son
significativos al 5\%. En la matriz de distancia, el rezago del PBI inicial
alcanza un valor de 0.2323 con $p=0.0932$, mientras que en la matriz queen el
coeficiente es 0.1979 con $p=0.2704$. En consecuencia, la evidencia de
desbordamientos directos provenientes del nivel inicial de los departamentos
vecinos es limitada.

La comparación entre matrices muestra que la importancia de la dependencia
espacial es sensible a la definición de vecindad. La banda inversa no mejora el
AIC respecto del OLS, mientras que la contigüidad queen produce una mejora
considerable en el modelo económico. Por ello, la conclusión central se apoya
en la estabilidad del coeficiente de convergencia y no únicamente en la
selección de una especificación espacial particular.

\begin{table}[H]
\centering
\caption{Resultados de los modelos globales y espaciales}
\label{tab:modelos_regresion}
\scriptsize
\resizebox{\textwidth}{!}{
\begin{tabular}{llrrrrrr}
\toprule
\textbf{Variable} &
\textbf{Modelo} &
\textbf{$\beta$ inicial} &
\textbf{$p(\beta)$} &
\textbf{Rezago inicial} &
\textbf{$p$ rezago} &
\textbf{$\lambda$} &
\textbf{AIC} \\
\midrule
\multirow{5}{*}{Cambio IDH}
& OLS
& 0.1637 & 0.0355 & --- & --- & --- & $-106.4696$ \\

& SEM, banda inversa
& 0.1712 & 0.0123 & --- & --- & $-0.3166$ & $-105.2509$ \\

& Durbin-SEM, banda inversa
& 0.1698 & 0.0139 & 0.0317 & 0.8515 & $-0.3321$ & $-103.2841$ \\

& SEM, queen
& 0.1882 & 0.0063 & --- & --- & $-0.3117$ & $-105.3478$ \\

& Durbin-SEM, queen
& 0.1782 & 0.0064 & 0.2782 & 0.0710 & $-0.4470$ & $-106.2370$ \\
\midrule
\multirow{5}{*}{Cambio $\ln$ PBIpc}
& OLS
& $-0.3282$ & 0.0004 & --- & --- & --- & 2.3566 \\

& SEM, banda inversa
& $-0.3666$ & $<0.001$ & --- & --- & 0.3036 & 3.2253 \\

& Durbin-SEM, banda inversa
& $-0.3471$ & $<0.001$ & 0.2323 & 0.0932 & 0.1120 & 3.0383 \\

& SEM, queen
& $-0.3792$ & $<0.001$ & --- & --- & 0.6304 & $-1.3505$ \\

& Durbin-SEM, queen
& $-0.3439$ & $<0.001$ & 0.1979 & 0.2704 & 0.5564 & $-0.4129$ \\
\bottomrule
\end{tabular}
}
\begin{flushleft}
\footnotesize
\textit{Nota.} $\beta$ corresponde al coeficiente del nivel inicial.
El rezago inicial representa la variable explicativa espacialmente rezagada.
Los valores de $\lambda$ corresponden al parámetro del error espacial.
\end{flushleft}
\end{table}

\subsection{Heterogeneidad espacial de los coeficientes}

La regresión geográficamente ponderada permite evaluar si la relación entre el
nivel inicial y el cambio acumulado varía entre departamentos. Los diagnósticos
se presentan en la Tabla \ref{tab:diagnostico_gwr}.

Para el IDH, el ancho de banda seleccionado es de aproximadamente 1919.1 km,
equivalente a más del 100\% de la distancia máxima considerada en el
procedimiento. Asimismo, los coeficientes locales se encuentran entre 0.1625 y
0.1652, con un rango de apenas 0.0027. Estos resultados indican que la GWR del
IDH se comporta prácticamente como una regresión global. No se detecta una
heterogeneidad territorial relevante en la pendiente estimada.

En el PBI per cápita relativo, el ancho de banda disminuye a 360.5 km, lo que
representa aproximadamente 19.1\% de la distancia máxima. Los coeficientes
locales varían entre $-0.4857$ y $-0.2771$, con un rango de 0.2086. Todos los
coeficientes conservan signo negativo, pero la intensidad de la convergencia
cambia entre departamentos.

La existencia de coeficientes económicos negativos en todo el territorio
confirma que la convergencia no depende únicamente de un reducido grupo de
departamentos. Sin embargo, la magnitud más intensa en determinadas zonas
sugiere que las condiciones productivas, la accesibilidad, la estructura
económica y la articulación regional influyen en la velocidad con la que los
territorios modifican su posición relativa.

\begin{table}[H]
\centering
\caption{Diagnóstico de las regresiones geográficamente ponderadas}
\label{tab:diagnostico_gwr}
\small
\begin{tabular}{lrrrrr}
\toprule
\textbf{Variable} &
\textbf{Banda, km} &
\textbf{Proporción máxima} &
\textbf{$\beta$ global} &
\textbf{$\beta$ local mínimo} &
\textbf{$\beta$ local máximo} \\
\midrule
IDH
& 1919.1 & 1.0142 & 0.1637 & 0.1625 & 0.1652 \\

PBI per cápita
& 360.5 & 0.1905 & $-0.3282$ & $-0.4857$ & $-0.2771$ \\
\bottomrule
\end{tabular}
\begin{flushleft}
\footnotesize
\textit{Nota.} La proporción máxima relaciona el ancho de banda seleccionado
con la distancia máxima entre los centroides departamentales.
\end{flushleft}
\end{table}

La Figura \ref{fig:gwr_paneles} representa los coeficientes locales y los
niveles de ajuste de las estimaciones. En el IDH se observa una superficie
prácticamente uniforme, coherente con el reducido rango de las pendientes
locales. En el PBI per cápita aparecen diferencias territoriales más visibles,
aunque todos los coeficientes mantienen el signo compatible con convergencia.

\begin{figure}[H]
\centering
\includegraphics[width=0.96\textwidth]{figura_4_gwr_paneles.png}
\caption{Resultados territoriales de las regresiones geográficamente ponderadas}
\label{fig:gwr_paneles}
\begin{flushleft}
\footnotesize
\textit{Nota.} Los paneles muestran la distribución departamental del
coeficiente local y del ajuste de las regresiones GWR para el IDH y el PBI
real per cápita relativo.
\end{flushleft}
\end{figure}

Debido al reducido número de observaciones, estos resultados deben
interpretarse de manera exploratoria. Un ancho de banda amplio, como el
obtenido para el IDH, puede reflejar una relación verdaderamente global, pero
también puede estar asociado con una potencia estadística limitada para
detectar variaciones locales. Por esta razón, la escasa heterogeneidad
estimada no debe interpretarse como prueba definitiva de homogeneidad
territorial.

\subsection{Relaciones no paramétricas}

La Figura \ref{fig:no_parametricas} compara la regresión lineal con los
suavizamientos LOESS y kernel. En el IDH, las tres aproximaciones mantienen una
relación predominantemente positiva entre el nivel inicial y el cambio
acumulado. La coincidencia general de las curvas respalda la conclusión de
ausencia de convergencia beta.

En el PBI per cápita relativo, las estimaciones muestran una relación negativa.
Los departamentos con una posición inicial inferior presentan, en promedio,
una variación posterior mayor. Aunque las curvas no paramétricas revelan
cambios moderados en la pendiente, no alteran el signo central identificado por
los modelos globales.

La comparación entre ambos indicadores demuestra que la convergencia regional
no es un proceso uniforme entre dimensiones. La aproximación relativa del
ingreso departamental coexistió con una dinámica distinta en desarrollo
humano. Esto implica que los mecanismos que favorecen la reducción de las
brechas económicas no necesariamente producen resultados equivalentes en
educación, salud y capacidades humanas.

\begin{figure}[H]
\centering
\includegraphics[width=0.96\textwidth]{figura_6_regresiones_no_parametricas.png}
\caption{Regresiones lineales y no paramétricas de convergencia}
\label{fig:no_parametricas}
\begin{flushleft}
\footnotesize
\textit{Nota.} Se comparan la recta OLS, el suavizamiento LOESS y la regresión
kernel para el IDH y el PBI real per cápita relativo.
\end{flushleft}
\end{figure}
